% =====================================================================
%  5. TIME MANAGEMENT
% =====================================================================
\section{Time Management}\label{sec:time}

% Owner colours used in the Gantt chart and its legend
\definecolor{ownerTKN}{HTML}{D50032}
\definecolor{ownerNMV}{HTML}{0F766E}
\definecolor{ownerHDN}{HTML}{D97706}
\definecolor{ownerNKNT}{HTML}{6D28D9}

The project runs from Sunday 20 September, when the team held its kick-off meeting, to Tuesday 8 December 2026. Its schedule is structured around the three \unitcode{} assessment deadlines, which divide the work into three stages, followed by a short closure period. Throughout this section, project weeks run from Monday to Sunday. Week~1 (W1) begins on Monday 5 October, and the planning weeks before it are numbered W$-$2 to W0.

\subsection{Project Stages and Assessment Deadlines}\label{sec:stages}

Table~\ref{tab:stages} maps each stage to its assessment, required submission and work packages. Each assessment submission is a hard deadline and is therefore treated as a fixed milestone in the schedule.

\begin{table}[htbp]
\centering
\caption{Project stages, assessment deadlines and submissions}
\label{tab:stages}
\begin{tabularx}{\textwidth}{@{}>{\raggedright\arraybackslash}p{0.2\textwidth} >{\raggedright\arraybackslash}p{0.19\textwidth} >{\raggedright\arraybackslash}X >{\raggedright\arraybackslash}p{0.12\textwidth}@{}}
  \toprule
  \textbf{Stage} & \textbf{Period and Deadline} & \textbf{Required Submission} & \textbf{WBS} \\
  \midrule
  \textbf{1. Planning and design} & W$-$2 to W1 (20 Sep -- 11 Oct)\newline \textbf{A1:} Sun 11 Oct & This Project Management Plan, including the initial front-end prototype & 1.1, 5.1, 5.2 \\
  \addlinespace
  \textbf{2. AI model development} (answers RQ1) & W2--W5 (12 Oct -- 2 Nov)\newline \textbf{A2:} Mon 2 Nov & Report (up to 3,000 words, PDF) and a ZIP file containing the processed datasets (CSV/JSON), a README explaining how to run the code, and the source code including the final models & 2, 3, 4.1--4.3, 4.5, 7.1 \\
  \addlinespace
  \textbf{3. Application development and evaluation} (answers RQ2--RQ5) & W5--W9 (2 Nov -- 6 Dec)\newline \textbf{A3:} Sun 6 Dec & Report (PDF), a ZIP file containing the React front end, the FastAPI back end with the integrated model and a README, and a demonstration video of up to 7 minutes & 4.4, 5.3--5.6, 6, 7.2, 7.3 \\
  \addlinespace
  \textbf{4. Closure} & W10 (7--8 Dec) & No submission; handover and lessons learned & 7.4 \\
  \bottomrule
\end{tabularx}
\end{table}

\subsection{Milestones}\label{sec:milestones}

Table~\ref{tab:milestones} lists the project's milestones. Assessment deadlines are shown in bold.

\begin{table}[htbp]
\centering
\caption{Project milestones}
\label{tab:milestones}
\begin{tabularx}{\textwidth}{@{}l l >{\raggedright\arraybackslash}X >{\raggedright\arraybackslash}p{0.33\textwidth}@{}}
  \toprule
  \textbf{ID} & \textbf{Date} & \textbf{Milestone} & \textbf{Completion Evidence} \\
  \midrule
  M1  & Sun 20 Sep & Project kick-off & Topic, team roles and working approach agreed (Meeting~1) \\
  M2  & Mon 5 Oct  & Dataset request sent & Email sent to the PRISMO dataset authors; fallback dataset downloaded \\
  M3  & Fri 9 Oct  & Ethics requirements confirmed & Tutor's advice on ethics approval for the user study recorded \\
  \textbf{M4}  & \textbf{Sun 11 Oct} & \textbf{Assignment 1 submitted} & Project Management Plan submitted \\
  M5  & Mon 12 Oct & Dataset decision & PRISMO dataset received, or switch to the fallback dataset confirmed \\
  M6  & Sun 25 Oct & Front-end prototype finalised & Prototype reviewed and approved by the team \\
  M7  & Tue 27 Oct & Feature selection complete & Subsets FS1--FS4 produced and documented \\
  M8  & Sat 31 Oct & Final model exported & Selected model and pipeline saved and loadable \\
  \textbf{M9}  & \textbf{Mon 2 Nov}  & \textbf{Assignment 2 submitted} & Model report and ZIP package submitted \\
  M10 & Sun 22 Nov & Core application integrated & All \textbf{Must} features working end to end \\
  M11 & Mon 30 Nov & User study complete & Sessions run and results analysed \\
  M12 & Wed 2 Dec  & Code freeze & No new features; only bug fixes and documentation \\
  \textbf{M13} & \textbf{Sun 6 Dec}  & \textbf{Assignment 3 submitted} & Report, ZIP package and demonstration video submitted \\
  M14 & Tue 8 Dec  & Project closed & Retrospective held and lessons learned recorded \\
  \bottomrule
\end{tabularx}
\end{table}

\subsection{Development Schedule}\label{sec:schedule}

Table~\ref{tab:schedule} shows the start and end dates, duration, owner and predecessors of every work package in the WBS (Section~\ref{sec:wbs}). The schedule was built around four principles:

\begin{enumerate}[nosep]
  \item \textbf{Work in parallel from Week 1.} No member waits on another stage to begin. While the dataset is being requested, the classifiers and the GA, PSO and SA algorithms are implemented and tested on sample data. Meanwhile, the Software Engineering member designs the architecture and prototype, and builds the pre-processing pipeline, which is later reused by the back end.
  \item \textbf{Make the dataset decision quickly.} Assignment~2 leaves only three weeks after the dataset is received, so the decision point (M5) is set one week after the request rather than two. The fallback dataset is downloaded in Week~1 so that a switch causes no further waiting.
  \item \textbf{Protect the assessment deadlines.} A code freeze (M12) four days before Assignment~3 leaves time for the demonstration video, final testing and the report.
  \item \textbf{Start the next stage early.} Work for Stage~3 that does not depend on the final model, such as the API skeleton and the prototype, begins during Stage~2.
\end{enumerate}

\paragraph{Critical path.} In Stage~2, the critical path runs through dataset acquisition, pre-processing, feature selection, priority-based feature formation, cross-validation, model comparison and the Assignment~2 report (2.1 $\rightarrow$ 2.3 $\rightarrow$ 3.2--3.4 $\rightarrow$ 3.5 $\rightarrow$ 4.2 $\rightarrow$ 4.3 $\rightarrow$ 7.1). It has almost no float, so any delay will be absorbed by contingency hours (Section~\ref{sec:wbs-dictionary}) or by reducing the number of experiment runs rather than by moving the deadline. In Stage~3, the critical path runs from model export through the API, integration, the user study and the Assignment~3 report (4.5 $\rightarrow$ 5.3 $\rightarrow$ 6.1 $\rightarrow$ 6.4 $\rightarrow$ 7.2), protected by the four-day buffer after the code freeze.

% =====================================================================
%  Development schedule 
% =====================================================================
{\renewcommand{\arraystretch}{1.15}\small
\begin{longtable}{@{}l >{\raggedright\arraybackslash}p{0.31\textwidth} l l r c >{\raggedright\arraybackslash}p{0.13\textwidth}@{}}
  \caption{Development schedule}\label{tab:schedule} \\
  \toprule
  \textbf{ID} & \textbf{Activity} & \textbf{Start} & \textbf{End} & \textbf{Days} & \textbf{Owner} & \textbf{Predecessors} \\
  \midrule
  \endfirsthead
  \multicolumn{7}{@{}l}{\small\itshape Table~\thetable{} continued} \\
  \toprule
  \textbf{ID} & \textbf{Activity} & \textbf{Start} & \textbf{End} & \textbf{Days} & \textbf{Owner} & \textbf{Predecessors} \\
  \midrule
  \endhead
  \midrule
  \multicolumn{7}{@{}r}{\small\itshape Continued on next page} \\
  \endfoot
  \bottomrule
  \multicolumn{7}{@{}l}{\footnotesize\parbox{\textwidth}{\vspace{0.3em}
  $^{a}$Implementation starts on sample data before the dataset arrives; final runs require the prepared dataset.
  $^{b}$An initial prototype is included in this plan (Section~\ref{sec:prototype}); it is refined until Week~3.
  $^{c}$Built against a mock model first; the real model is connected after M8.
  $^{d}$Unit tests start earlier; system tests begin after integration.
  Days are calendar days, inclusive. Where two owners are listed, the first is accountable and the second shares the work.}} \\
  \endlastfoot
  \multicolumn{7}{@{}l}{\cellcolor{brandprimary!8}\textbf{\sffamily 1 \quad Project Management}} \\
  1.1 & Project initiation \& planning & Sun 20 Sep & Sun 11 Oct & 22 & TKN & -- \\
  1.2 & Meetings \& progress reporting & Sun 20 Sep & Sun 6 Dec & 78 & TKN & -- \\
  1.3 & Risk \& change management & Mon 5 Oct & Sun 6 Dec & 63 & TKN & 1.1 \\
  \multicolumn{7}{@{}l}{\cellcolor{brandprimary!8}\textbf{\sffamily 2 \quad Data Collection \& Preparation}} \\
  2.1 & Dataset request \& acquisition & Mon 5 Oct & Mon 12 Oct & 8 & TKN & -- \\
  2.2 & Exploratory data analysis & Tue 13 Oct & Sun 18 Oct & 6 & NMV & 2.1 \\
  2.3 & Data cleaning \& pre-processing & Tue 13 Oct & Wed 21 Oct & 9 & HDN & 2.1 \\
  2.4 & Feature preparation & Tue 13 Oct & Tue 20 Oct & 8 & TKN & 2.1 \\
  \multicolumn{7}{@{}l}{\cellcolor{brandprimary!8}\textbf{\sffamily 3 \quad Feature Subset Selection}} \\
  3.1 & Baseline (full feature set) & Thu 22 Oct & Fri 23 Oct & 2 & NKNT & 2.3, 2.4, 4.1 \\
  3.2 & GA feature selection & Mon 12 Oct & Sun 25 Oct & 14 & NKNT & 2.3, 2.4$^{a}$ \\
  3.3 & PSO feature selection & Mon 12 Oct & Sun 25 Oct & 14 & NMV & 2.3, 2.4$^{a}$ \\
  3.4 & SA feature selection & Mon 12 Oct & Sun 25 Oct & 14 & TKN & 2.3, 2.4$^{a}$ \\
  3.5 & Priority-based feature formation & Mon 26 Oct & Tue 27 Oct & 2 & NMV & 3.2--3.4 \\
  \multicolumn{7}{@{}l}{\cellcolor{brandprimary!8}\textbf{\sffamily 4 \quad Classification \& Evaluation}} \\
  4.1 & Classifier implementation & Mon 5 Oct & Sun 18 Oct & 14 & NKNT, TKN & --$^{a}$ \\
  4.2 & Cross-validation experiments & Tue 27 Oct & Thu 29 Oct & 3 & NKNT & 3.1, 3.5 \\
  4.3 & Comparison \& model selection & Thu 29 Oct & Fri 30 Oct & 2 & NMV & 4.2 \\
  4.4 & Prediction explanations & Mon 2 Nov & Sun 8 Nov & 7 & NKNT & 4.5 \\
  4.5 & Model export & Fri 30 Oct & Sat 31 Oct & 2 & NKNT & 4.3 \\
  \multicolumn{7}{@{}l}{\cellcolor{brandprimary!8}\textbf{\sffamily 5 \quad Application Development}} \\
  5.1 & System architecture design & Mon 5 Oct & Fri 9 Oct & 5 & HDN & -- \\
  5.2 & Front-end prototype & Mon 5 Oct & Sun 25 Oct & 21 & HDN & 5.1$^{b}$ \\
  5.3 & Prediction API (FastAPI) & Mon 26 Oct & Sun 15 Nov & 21 & HDN & 5.1, 4.5$^{c}$ \\
  5.4 & Web application (React) & Mon 2 Nov & Sun 22 Nov & 21 & HDN & 5.2 \\
  5.5 & Decision support \& study mode & Mon 9 Nov & Fri 20 Nov & 12 & NKNT, TKN & 4.4, 5.3 \\
  5.6 & Interactive visualisations & Mon 2 Nov & Fri 20 Nov & 19 & NMV, TKN & 4.3 \\
  \multicolumn{7}{@{}l}{\cellcolor{brandprimary!8}\textbf{\sffamily 6 \quad Integration, Testing \& User Evaluation}} \\
  6.1 & System integration & Mon 16 Nov & Sun 22 Nov & 7 & HDN & 5.3--5.6 \\
  6.2 & Functional \& performance testing & Mon 16 Nov & Wed 2 Dec & 17 & NKNT & 6.1$^{d}$ \\
  6.3 & User study design \& ethics & Mon 2 Nov & Fri 20 Nov & 19 & NMV & M3 \\
  6.4 & User study sessions \& analysis & Mon 23 Nov & Mon 30 Nov & 8 & NMV, TKN & 6.1, 6.3 \\
  \multicolumn{7}{@{}l}{\cellcolor{brandprimary!8}\textbf{\sffamily 7 \quad Closure \& Submissions}} \\
  7.1 & Assignment 2 report \& code package & Mon 26 Oct & Mon 2 Nov & 8 & TKN & 4.3 \\
  7.2 & Assignment 3 report & Mon 23 Nov & Sun 6 Dec & 14 & TKN & 6.1 \\
  7.3 & Demonstration video & Wed 2 Dec & Sat 5 Dec & 4 & HDN, TKN & M12 \\
  7.4 & Handover \& lessons learned & Mon 7 Dec & Tue 8 Dec & 2 & TKN & M13 \\
\end{longtable}
}


\subsection{Gantt Chart}\label{sec:gantt}

Figure~\ref{fig:gantt} presents the schedule as a Gantt chart. Bars are coloured by owner, group bars summarise each WBS branch, and diamonds mark milestones and assessment deadlines.

\begin{figure}[p]
\centering
\resizebox{\textwidth}{!}{% =====================================================================
%  Gantt chart =====================================================================
\begin{ganttchart}[
    time slot format=isodate,
    x unit=0.13cm,
    y unit title=0.45cm,
    y unit chart=0.4cm,
    title height=1,
    title/.append style={draw=brandprimary!40, fill=brandprimary!8},
    title label font=\tiny\sffamily,
    hgrid={draw=black!6},
    vgrid={*6{draw=none}, {draw=black!20}},
    bar height=0.55,
    bar top shift=0.225,
    bar/.append style={draw=none},
    bar label font=\scriptsize\sffamily,
    group/.append style={fill=brandprimary, draw=brandprimary},
    group label font=\scriptsize\sffamily\bfseries,
    group height=0.25,
    group top shift=0.375,
    group peaks height=0.15,
    milestone/.append style={fill=brandaccent, draw=brandaccent},
    milestone label font=\scriptsize\sffamily\itshape,
    milestone height=0.6,
    milestone top shift=0.2
  ]{2026-09-14}{2026-12-13}
  \gantttitlecalendar{month=shortname} \\
  \gantttitle{W$-$2}{7} \gantttitle{W$-$1}{7} \gantttitle{W0}{7} \gantttitle{W1}{7} \gantttitle{W2}{7} \gantttitle{W3}{7} \gantttitle{W4}{7} \gantttitle{W5}{7} \gantttitle{W6}{7} \gantttitle{W7}{7} \gantttitle{W8}{7} \gantttitle{W9}{7} \gantttitle{W10}{7} \\
  \ganttmilestone{A1 -- Assignment 1 due (this plan)}{2026-10-11} \\
  \ganttmilestone{A2 -- Assignment 2 due (AI model)}{2026-11-02} \\
  \ganttmilestone{A3 -- Assignment 3 due (web app)}{2026-12-06} \\
  \ganttgroup{1 Project Management}{2026-09-20}{2026-12-06} \\
  \ganttbar[bar/.append style={fill=ownerTKN}]{1.1 Project initiation \& planning}{2026-09-20}{2026-10-11} \\
  \ganttbar[bar/.append style={fill=ownerTKN}]{1.2 Meetings \& progress reporting}{2026-09-20}{2026-12-06} \\
  \ganttbar[bar/.append style={fill=ownerTKN}]{1.3 Risk \& change management}{2026-10-05}{2026-12-06} \\
  \ganttmilestone{M1 Project kick-off}{2026-09-20} \\
  \ganttgroup{2 Data Collection \& Preparation}{2026-10-05}{2026-10-21} \\
  \ganttbar[bar/.append style={fill=ownerTKN}]{2.1 Dataset request \& acquisition}{2026-10-05}{2026-10-12} \\
  \ganttbar[bar/.append style={fill=ownerNMV}]{2.2 Exploratory data analysis}{2026-10-13}{2026-10-18} \\
  \ganttbar[bar/.append style={fill=ownerHDN}]{2.3 Data cleaning \& pre-processing}{2026-10-13}{2026-10-21} \\
  \ganttbar[bar/.append style={fill=ownerTKN}]{2.4 Feature preparation}{2026-10-13}{2026-10-20} \\
  \ganttmilestone{M5 Dataset decision}{2026-10-12} \\
  \ganttgroup{3 Feature Subset Selection}{2026-10-12}{2026-10-27} \\
  \ganttbar[bar/.append style={fill=ownerNKNT}]{3.1 Baseline (full feature set)}{2026-10-22}{2026-10-23} \\
  \ganttbar[bar/.append style={fill=ownerNKNT}]{3.2 GA feature selection}{2026-10-12}{2026-10-25} \\
  \ganttbar[bar/.append style={fill=ownerNMV}]{3.3 PSO feature selection}{2026-10-12}{2026-10-25} \\
  \ganttbar[bar/.append style={fill=ownerTKN}]{3.4 SA feature selection}{2026-10-12}{2026-10-25} \\
  \ganttbar[bar/.append style={fill=ownerNMV}]{3.5 Priority-based feature formation}{2026-10-26}{2026-10-27} \\
  \ganttmilestone{M7 Feature selection complete}{2026-10-27} \\
  \ganttgroup{4 Classification \& Evaluation}{2026-10-05}{2026-11-08} \\
  \ganttbar[bar/.append style={fill=ownerNKNT}, bar height=0.275, bar top shift=0.225]{4.1 Classifier implementation}{2026-10-05}{2026-10-18}
  \ganttbar[bar/.append style={fill=ownerTKN}, bar height=0.275, bar top shift=0.5]{}{2026-10-05}{2026-10-18} \\
  \ganttbar[bar/.append style={fill=ownerNKNT}]{4.2 Cross-validation experiments}{2026-10-27}{2026-10-29} \\
  \ganttbar[bar/.append style={fill=ownerNMV}]{4.3 Comparison \& model selection}{2026-10-29}{2026-10-30} \\
  \ganttbar[bar/.append style={fill=ownerNKNT}]{4.4 Prediction explanations}{2026-11-02}{2026-11-08} \\
  \ganttbar[bar/.append style={fill=ownerNKNT}]{4.5 Model export}{2026-10-30}{2026-10-31} \\
  \ganttmilestone{M8 Final model exported}{2026-10-31} \\
  \ganttgroup{5 Application Development}{2026-10-05}{2026-11-22} \\
  \ganttbar[bar/.append style={fill=ownerHDN}]{5.1 System architecture design}{2026-10-05}{2026-10-09} \\
  \ganttbar[bar/.append style={fill=ownerHDN}]{5.2 Front-end prototype}{2026-10-05}{2026-10-25} \\
  \ganttbar[bar/.append style={fill=ownerHDN}]{5.3 Prediction API (FastAPI)}{2026-10-26}{2026-11-15} \\
  \ganttbar[bar/.append style={fill=ownerHDN}]{5.4 Web application (React)}{2026-11-02}{2026-11-22} \\
  \ganttbar[bar/.append style={fill=ownerNKNT}, bar height=0.275, bar top shift=0.225]{5.5 Decision support \& study mode}{2026-11-09}{2026-11-20}
  \ganttbar[bar/.append style={fill=ownerTKN}, bar height=0.275, bar top shift=0.5]{}{2026-11-09}{2026-11-20} \\
  \ganttbar[bar/.append style={fill=ownerNMV}, bar height=0.275, bar top shift=0.225]{5.6 Interactive visualisations}{2026-11-02}{2026-11-20}
  \ganttbar[bar/.append style={fill=ownerTKN}, bar height=0.275, bar top shift=0.5]{}{2026-11-02}{2026-11-20} \\
  \ganttmilestone{M6 Prototype finalised}{2026-10-25} \\
  \ganttgroup{6 Integration, Testing \& User Evaluation}{2026-11-02}{2026-12-02} \\
  \ganttbar[bar/.append style={fill=ownerHDN}]{6.1 System integration}{2026-11-16}{2026-11-22} \\
  \ganttbar[bar/.append style={fill=ownerNKNT}]{6.2 Functional \& performance testing}{2026-11-16}{2026-12-02} \\
  \ganttbar[bar/.append style={fill=ownerNMV}]{6.3 User study design \& ethics}{2026-11-02}{2026-11-20} \\
  \ganttbar[bar/.append style={fill=ownerNMV}, bar height=0.275, bar top shift=0.225]{6.4 User study sessions \& analysis}{2026-11-23}{2026-11-30}
  \ganttbar[bar/.append style={fill=ownerTKN}, bar height=0.275, bar top shift=0.5]{}{2026-11-23}{2026-11-30} \\
  \ganttmilestone{M10 Core application integrated}{2026-11-22} \\
  \ganttmilestone{M11 User study complete}{2026-11-30} \\
  \ganttgroup{7 Closure \& Submissions}{2026-10-26}{2026-12-08} \\
  \ganttbar[bar/.append style={fill=ownerTKN}]{7.1 Assignment 2 report \& code package}{2026-10-26}{2026-11-02} \\
  \ganttbar[bar/.append style={fill=ownerTKN}]{7.2 Assignment 3 report}{2026-11-23}{2026-12-06} \\
  \ganttbar[bar/.append style={fill=ownerHDN}, bar height=0.275, bar top shift=0.225]{7.3 Demonstration video}{2026-12-02}{2026-12-05}
  \ganttbar[bar/.append style={fill=ownerTKN}, bar height=0.275, bar top shift=0.5]{}{2026-12-02}{2026-12-05} \\
  \ganttbar[bar/.append style={fill=ownerTKN}]{7.4 Handover \& lessons learned}{2026-12-07}{2026-12-08} \\
  \ganttmilestone{M12 Code freeze}{2026-12-02}
\end{ganttchart}
}

\smallskip
{\scriptsize\sffamily
\textcolor{ownerTKN}{\rule{9pt}{6pt}}~TKN (\memberOne) \quad
\textcolor{ownerNMV}{\rule{9pt}{6pt}}~NMV (\memberTwo) \quad
\textcolor{ownerHDN}{\rule{9pt}{6pt}}~HDN (\memberThree) \quad
\textcolor{ownerNKNT}{\rule{9pt}{6pt}}~NKNT (\memberFour) \quad
\textcolor{brandaccent}{$\blacklozenge$}~Milestone}
\caption{Gantt chart for the \projectname{} project (W1 begins Monday 5 October 2026)}
\label{fig:gantt}
\end{figure}
